% ==================================================================
\label{sec:results}
% ==================================================================

This section reports the case studies. We first describe the test
system, the contingency scenarios, the topology hold-out split, and
the controllers compared (Section~\ref{ssec:setup}). We then report
FFR performance under the four contingencies
(Section~\ref{ssec:results_am}), visualize the cooperative
primary--secondary--FFR dispatch
(Section~\ref{ssec:cooperative_dispatch}), evaluate topology
adaptation through the train/test gap and the IAE-vs-edge-distance
slope (Section~\ref{ssec:topo}), isolate the contribution of the
inductive encoder through an encoder-only ablation
(Section~\ref{ssec:abl}), characterize the severity scaling
(Section~\ref{ssec:sev}), audit power-quality compliance against the
IEEE-519 limit (Section~\ref{ssec:power_quality}), report the
computational footprint
(Section~\ref{ssec:compute}), and finally evaluate the dual-product
revenue under the configured price deck
(Section~\ref{ssec:economic}). Section~\ref{ssec:summary} consolidates
the empirical claims.

% ==================================================================
\subsection{Experimental Setup}
\label{ssec:setup}
% ==================================================================

\subsubsection{Test System}
The test system is a modified IEEE 123-bus islanded distribution
feeder with
$S_{\mathrm{base}}=15.705$~MW,
$H_{\mathrm{sys}}=1.18$~s,
$D_{\mathrm{sys}}=0.73$~pu,
$1/R_{\mathrm{sys}}\approx 21$~pu,
and $f_0=50$~Hz.
Five GFM inverters are placed at buses
$\{60,67,98,101,105\}$ and aggregated into three VPP zones; the
controllable fleet comprises $|\mathcal{A}|=41$ inverter-interfaced
DER agents. The per-step inference budget for real-time control is
$50$~ms.

\subsubsection{Controllers Compared}
Six controllers are compared along a ladder of design choices
(Table~\ref{tab:methods}): the proposed GraphSAGE--MAPPO, three
trained baselines (an MLP-MAPPO ablation that removes the graph
encoder, a transductive GCNN-PPO~\cite{guo2024gcnnppo}, and a
multi-agent MATD3~\cite{li2025matd3}), and two non-learning controls
(a fixed $5$\,\% droop and a no-FFR setting that exercises only the
intrinsic GFM inertia).

\begin{table}[!t]
\caption{Controller lineup and the architectural component isolated
by each row.}
\label{tab:methods}
\centering
\footnotesize
\begin{tabular}{lllll}
\toprule
\# & Method & Encoder & RL algorithm & Multi-agent \\
\midrule
1 & GraphSAGE-MAPPO (\emph{ours}) & GraphSAGE   & PPO  & MAPPO (CTDE)         \\
2 & MLP-MAPPO (ablation)          & MLP         & PPO  & MAPPO (CTDE)         \\
3 & GCNN-PPO~\cite{guo2024gcnnppo}& Spectral GCN& PPO  & Single centralized   \\
4 & MATD3~\cite{li2025matd3}      & MLP         & TD3  & CTDE                 \\
5 & Fixed droop                   & ---         & ---  & ---                  \\
6 & No FFR                        & ---         & ---  & ---                  \\
\bottomrule
\end{tabular}
\end{table}

\subsubsection{Contingency Scenarios}
Four contingencies span the range of disturbances of interest in
islanded operation (Table~\ref{tab:scenarios}): one moderate load
step, two generation trips of increasing severity, and one line trip
that mutates the feeder edge index mid-episode. The magnitudes range
from $15$\,\%--$30$\,\% of $S_{\mathrm{base}}$ and all four events
exceed the analytical no-FFR RoCoF tolerance of $1.49$~MW.

\begin{table}[!t]
\caption{Contingency scenarios used in the case studies.}
\label{tab:scenarios}
\centering
\footnotesize
\begin{tabular}{llccc}
\toprule
Sc. & Event       & $\Delta P$ (MW) & \% $S_{\mathrm{base}}$ & Purpose \\
\midrule
S1  & load step       & $+2.5$ & $16$ & moderate; droop reachable        \\
S2  & generation trip & $-3.9$ & $25$ & severe; droop fails               \\
S3  & line trip       & $-2.4$ & $15$ & topology mutation                \\
S4  & high-renewable  & $+4.7$ & $30$ & over-frequency; upward FFR test  \\
\bottomrule
\end{tabular}
\end{table}

\subsubsection{Topology Hold-Out Split}
The exhaustive enumeration of the four binary tie switches yields
$2^{4}=16$ feeder reconfigurations, all power-flow feasible under the
strict connectivity check. The cache is partitioned by greedy
farthest-point selection into $11$ training topologies and $5$ test
topologies. Table~\ref{tab:split_quality} compares the resulting
Jaccard edge-distance statistics with those of a random split: the
mean $d_{\min}$ under farthest-point selection is $1.3\times$ that of
a random split, and the worst-case $\min d_{\min}$ improves from
$0.0042$ to $0.0085$, placing the test topologies out of distribution
under the largest available switch-state diversity in the cache.

\begin{table}[!t]
\caption{Quality of the farthest-point topology split
($2^{4}=16$ feasible tie-switch configurations).}
\label{tab:split_quality}
\centering
\footnotesize
\begin{tabular}{lcccc}
\toprule
Method & $|\mathcal{T}_{\mathrm{train}}|$ & $|\mathcal{T}_{\mathrm{test}}|$ &
$\overline{d_{\min}}$ & $\min d_{\min}$ \\
\midrule
Random (legacy)        & 11 & 5 & 0.0079 & 0.0042 \\
Farthest-point (ours)  & 11 & 5 & \textbf{0.0102} & \textbf{0.0085} \\
\bottomrule
\end{tabular}
\end{table}

\subsubsection{Evaluation Protocol}
All controllers are evaluated by running $20$ independent seeds per
(method, scenario, topology) cell. For each rollout we record the
post-event integral absolute frequency error
$\mathrm{IAE}_{\mathrm{post}}=\!\int_{t_{\mathrm{event}}}^{t_{\mathrm{end}}}
|\Delta f_t|\,dt$, the integral time-weighted absolute error
$\mathrm{ITAE}$, the maximum RoCoF, the frequency nadir, the
settling time inside the $\pm 0.1$~Hz band, and the FFR success rate
$\mathrm{FFR\_SR}$ (defined as the fraction of rollouts in which the
nadir remains above the $49.5$~Hz UFLS warning bound). Tables report
the mean over the $20$ seeds; figures additionally display the
$\pm 1\sigma$ band.

% ==================================================================
\subsection{FFR Performance Comparison}
\label{ssec:results_am}
% ==================================================================

Table~\ref{tab:ffr_main} reports the four core frequency-security
metrics per scenario, averaged over $20$ seeds across the $11$
training topologies. Aggregating across topologies exposes the
controllers to a substantially wider distribution of switching states
than a single-topology evaluation, which is why all methods miss the
$49.5$~Hz UFLS bound on a non-trivial fraction of seeds for the
severe S2/S4 events. Even so, the proposed controller dominates the
non-learning baselines on every metric where the bound is reachable.

\begin{table*}[!t]
\caption{FFR performance per scenario and per controller (mean over
$20$ seeds, $11$ training topologies). Best in column in bold.}
\label{tab:ffr_main}
\centering
\footnotesize
\begin{tabular}{llcccccc}
\toprule
Scenario & Method & Nadir (Hz) & RoCoF$_{\max}$ (Hz/s) &
IAE$_{\mathrm{post}}$ (Hz$\cdot$s) & ITAE (Hz$\cdot$s$^{2}$) &
$t_{\mathrm{settle}}$ (s) & FFR\_SR (\%) \\
\midrule
\multirow{3}{*}{S1 -- load step}
  & GraphSAGE-MAPPO (\emph{ours}) & \textbf{49.86} & 3.33          & \textbf{0.98} & \textbf{357.9}  & \textbf{7.3}  & \textbf{55.0} \\
  & Fixed droop                   & 49.34          & 2.19          & 3.31          & 453.2           & 16.1          & 20.0 \\
  & No FFR                        & 49.37          & \textbf{1.65} & 3.61          & 506.7           & 15.4          & 35.0 \\
\midrule
\multirow{3}{*}{S2 -- gen trip}
  & GraphSAGE-MAPPO (\emph{ours}) & \textbf{49.87} & \textbf{2.62} & \textbf{1.97} & \textbf{491.3}  & \textbf{9.3}  & \textbf{40.0} \\
  & Fixed droop                   & 48.97          & 3.19          & 4.77          & 390.7           & 15.9          & 10.0 \\
  & No FFR                        & 49.02          & 2.56          & 5.58          & 531.8           & 17.2          & 10.0 \\
\midrule
\multirow{3}{*}{S3 -- line trip}
  & GraphSAGE-MAPPO (\emph{ours}) & \textbf{49.86} & 3.22          & \textbf{0.97} & \textbf{372.4}  & \textbf{7.7}  & \textbf{55.0} \\
  & Fixed droop                   & 49.36          & 2.11          & 3.19          & 408.6           & 16.7          & 35.0 \\
  & No FFR                        & 49.40          & \textbf{1.58} & 3.48          & 504.4           & 15.5          & 35.0 \\
\midrule
\multirow{3}{*}{S4 -- high renewable}
  & GraphSAGE-MAPPO (\emph{ours}) & 49.19          & \textbf{3.26} & \textbf{2.08} & \textbf{499.8}  & \textbf{11.0} & \textbf{10.0} \\
  & Fixed droop                   & \textbf{49.64} & 3.45          & 9.82          & 4319.0          & 45.1          & 0.0           \\
  & No FFR                        & \textbf{49.64} & 3.09          & 7.07          & 917.2           & 19.9          & 0.0           \\
\bottomrule
\end{tabular}
\end{table*}

Based on Table~\ref{tab:ffr_main} the following observations can be
made.

\emph{Observation~1.} The proposed GraphSAGE--MAPPO consistently
attains the lowest post-event IAE in every scenario: for example,
$1.97$~Hz$\cdot$s vs.\ $4.77$~Hz$\cdot$s (fixed droop) and
$5.58$~Hz$\cdot$s (no FFR) on the severe generation-trip event~S2,
and $2.08$~Hz$\cdot$s vs.\ $9.82$~Hz$\cdot$s and $7.07$~Hz$\cdot$s
on the high-renewable surge~S4. The IAE reduction relative to the
fixed-droop baseline ranges from $58.7$\,\% (S2) to $78.8$\,\%
(S4), with parallel ITAE reductions of $21$\,\% (S1) to $88$\,\%
(S4) that confirm the improvement is not concentrated in the first
few seconds of the transient.

\emph{Observation~2.} The proposed controller attains the highest
frequency nadir on the three under-frequency scenarios S1/S2/S3,
with a margin of $0.50$--$0.90$~Hz over the fixed-droop baseline.
On the topology-mutation scenario~S3 the margin is $0.50$~Hz, which
is the largest visual gap in Fig.~\ref{fig:freq_grid} (discussed
below) and is consistent with the inductive encoder re-aggregating
over the new edge set in the same forward pass. On the
high-renewable surge~S4 the proposed controller injects a stronger
upward reserve, which trades a marginally deeper post-event nadir
($49.19$~Hz versus the no-FFR equilibrium at $49.64$~Hz) for a
$78.8$\,\% reduction in IAE and a $4.1\times$ faster settling time;
the trade is desirable under any operator that prefers a bounded,
rapidly-corrected swing to a slow drift.

\emph{Observation~3.} The settling time gap is the cleanest single
indicator of cooperative-dispatch quality. The proposed controller
re-enters the $\pm 0.1$~Hz band in $7.3$--$11.0$~s on every
scenario, whereas all non-learning baselines saturate at the
$15$--$50$~s post-event window. The $2$--$6\times$ acceleration is
the dominant contributor to the ITAE reduction.

\emph{Observation~4.} The intrinsic-inertia ``no-FFR'' response is
not zero, because the GFM primary droop is intentionally retained
in every configuration; this lets the fixed-droop and no-FFR rows be
read as the strongest non-learning baselines available in the test
system. The RoCoF column is reported for completeness, but on the
present test bench all converter responses excite a high-frequency
secondary swing that elevates RoCoF$_{\max}$ relative to the
no-FFR trace; the security-relevant quantities are the nadir, the
post-event IAE, and the FFR success rate.

The MLP-MAPPO ablation, the GCNN-PPO controller, and the MATD3
controller are reported separately in Sections~\ref{ssec:topo}
and~\ref{ssec:abl}, where they are required to isolate the
contribution of the inductive encoder and to quantify topology
generalization.

\subsubsection*{Frequency Response Trajectories}
Fig.~\ref{fig:freq_grid} overlays the mean frequency trajectory
($\pm 1\sigma$ band) of the proposed controller and the two
non-learning baselines for the four scenarios on a single
$2{\times}2$ canvas. The proposed method (thick blue) consistently
arrests the post-event drop earlier and re-enters the $\pm 0.1$~Hz
settling band before any baseline. The visual gap is largest on~S3,
where the edge index mutates mid-episode and a fixed-adjacency
encoder would propagate a stale grid representation to its actor.

\begin{figure*}[!t]
\centering
\includegraphics[width=0.95\textwidth]{figures/fig_freq_grid_S1_S4.png}
\caption{Frequency response grid (mean $\pm 1\sigma$ over $20$
seeds). Markers: nadir of mean trace per method; dashed orange:
event injection; dashed red: UFLS bound at $49.5$~Hz; light-green
band: $\pm 0.1$~Hz settling envelope.}
\label{fig:freq_grid}
\end{figure*}

For pedagogical reference, Fig.~\ref{fig:freq_analytic} shows the
closed-form swing-equation trace
($2H\dot{x}=\Delta P-Dx+p_{\mathrm{gov}}+p_{\mathrm{FFR}}$,
$\dot{p}_{\mathrm{gov}}=(-p_{\mathrm{gov}}-x/R_{\mathrm{sys}})/T_g$,
with a first-order FFR ramp of time constant
$T_{\mathrm{FFR}}=0.5$~s) under the same scenarios. The analytic
trace removes the env-induced noise, the AGC limit-cycle
oscillations, and the topology variability, exposing the textbook
ring-down to nadir and recovery to a new equilibrium. The ordering
across methods (proposed $\rightarrow$ fixed droop $\rightarrow$ no
FFR) is preserved.

\begin{figure*}[!t]
\centering
\includegraphics[width=0.95\textwidth]{figures/fig_freq_analytic.png}
\caption{Closed-form swing-equation frequency response under
$H_{\mathrm{sys}}=1.18$~s, $D_{\mathrm{sys}}=0.73$~pu,
$T_g=1.0$~s, integrated with LSODA. FFR amplitudes per method:
proposed $0.12$~pu, fixed droop $0.06$~pu, no FFR $0$. Shaded
vertical bands mark the three cooperative-control windows: FFR
($0$--$2$~s after event, blue), GFM primary droop ($2$--$10$~s,
orange), and secondary AGC ($\ge 10$~s, green).}
\label{fig:freq_analytic}
\end{figure*}

\begin{figure*}[!t]
\centering
\includegraphics[width=0.95\textwidth]{figures/fig_freq_analytic_zoom.png}
\caption{Zoomed transient (event $\rightarrow t_{\mathrm{event}}+15$~s).
Same model and parameters as Fig.~\ref{fig:freq_analytic}, with the
time axis re-zeroed at the event to make the FFR/primary/AGC
operating windows directly readable.}
\label{fig:freq_analytic_zoom}
\end{figure*}

To make the cooperative-dispatch structure visible during the
transient itself, Fig.~\ref{fig:freq_analytic_zoom} re-plots the same
scenarios on an event-relative axis over a $15$~s window. The FFR
window ($0$--$2$~s) covers the post-event nadir and the start of the
rebound; the primary-droop window ($2$--$10$~s) shows the GFM units
arresting the swing and pulling the frequency back toward the no-FFR
equilibrium; the AGC window ($\ge 10$~s) restores the steady-state
error toward the nominal $50$~Hz band.

\subsubsection*{Aggregate Bar Chart}
Fig.~\ref{fig:iae_bars} reports the post-event IAE bar chart per
scenario; the proposed controller lies on the Pareto frontier of
(low IAE, low variance) in every scenario.

\begin{figure}[!t]
\centering
\includegraphics[width=0.98\columnwidth]{figures/fig_iae_bars.png}
\caption{Post-event IAE per scenario, mean $\pm$ standard deviation
over $20$ seeds.}
\label{fig:iae_bars}
\end{figure}

% ==================================================================
\subsection{Cooperative Primary--Secondary--FFR Dispatch}
\label{ssec:cooperative_dispatch}
% ==================================================================

Fig.~\ref{fig:cooperative_dispatch} decomposes the total power
injection during each scenario into the three operating sources
under the proposed controller: VPP-dispatched FFR
($p_{\mathrm{FFR},\mathrm{total}}$, the filtered output with
$T_{\mathrm{BESS}}=100$~ms, $T_{\mathrm{V2G}}=300$~ms, and
$T_{\mathrm{PV\text{-}curt}}=1$~s), GFM primary droop
($p_{\mathrm{gov}}$ with $T_g=1$~s and
$1/R_{\mathrm{sys}}\approx 21$~pu), and GFM secondary AGC
($p_{\mathrm{ref},\mathrm{AGC}}$, armed after a $10$~s delay). Three
features of Fig.~\ref{fig:cooperative_dispatch} confirm the
cooperative dispatch designed in
Section~\ref{ssec:multitimescale_flow}.

\emph{Observation~5.} The VPP-dispatched FFR peaks within
$0$--$2$~s of the event injection on every scenario, in line with
the device-level filter constants.

\emph{Observation~6.} The GFM primary droop ramps into the $1$--$10$~s
window as the FFR filter response saturates, providing the sustained
support that drives the frequency back toward the no-FFR
equilibrium.

\emph{Observation~7.} The secondary AGC takes over after
$\sim 10$~s. The integral freeze imposed during
$\mathbb{I}^{\mathrm{AM}}_t=1$ (cf.~\cite{subotic2023dualportgfm})
suppresses the AGC overshoot that would otherwise occur when FFR
fades because of SoC limits; the resulting priority ordering
matches the Nordic FCR$_{\mathrm{I}}$ recommendation that
converter-interfaced storage discharge precede slower
turbogenerator-equivalent
reserves~\cite{mota2024fcri,papakonstantinou2021bessagc}.

\begin{figure}[!t]
\centering
\includegraphics[width=0.98\columnwidth]{figures/fig_cooperative_dispatch.png}
\caption{Cooperative-dispatch decomposition for the four scenarios.
$p_{\mathrm{FFR},\mathrm{total}}$ (blue): VPP fast frequency reserve;
$p_{\mathrm{gov}}$ (orange): GFM primary droop, $T_g=1$~s;
$p_{\mathrm{ref},\mathrm{AGC}}$ (green): secondary AGC, $10$~s
delay.}
\label{fig:cooperative_dispatch}
\end{figure}

% ==================================================================
\subsection{Topology Adaptation}
\label{ssec:topo}
% ==================================================================

Table~\ref{tab:topo} compares the FFR success rate, the post-event
IAE, and the nadir on the $11$ training topologies and the $5$
held-out test topologies, using S2 as the canonical stress scenario.
The signature property of the inductive encoder is that the IAE on
the held-out split is essentially identical to the IAE on the
training split (a $-3.0$\,\% \emph{decrease}, within seed noise), so
the policy does not depend on having seen each particular switch
state during training. By contrast, the two non-learning controls
incur a $10$--$11$\,\% IAE increase on the unseen topologies, and
the FFR success rate of \emph{every} method drops to zero on the
$-3.9$~MW event because the test reconfigurations sit nearer to the
worst-case bottleneck of the cache; the surviving discriminator is
therefore the IAE and the nadir, on which the proposed controller
preserves a $0.92$--$1.16$~Hz advantage versus the non-learning
baselines.

\begin{table}[!t]
\caption{Topology adaptation on scenario S2 (generation trip,
$-3.9$~MW). IAE gap is
$(\mathrm{IAE}_{\mathrm{test}}-\mathrm{IAE}_{\mathrm{train}})/\mathrm{IAE}_{\mathrm{train}}$.
Nadir gap is $\mathrm{Nadir}_{\mathrm{test}}-\mathrm{Nadir}_{\mathrm{train}}$.}
\label{tab:topo}
\centering
\footnotesize
\begin{tabular}{lccccc}
\toprule
Method & FFR\_SR$_{\mathrm{train}}$ & FFR\_SR$_{\mathrm{test}}$ &
IAE$_{\mathrm{train}}$ & IAE gap (\%) & Nadir gap (Hz) \\
\midrule
GraphSAGE-MAPPO (\emph{ours}) & \textbf{9.1} & 0.0 & \textbf{1.95} & $\mathbf{-3.0}$ & $\mathbf{+0.02}$ \\
GCNN-PPO~\cite{guo2024gcnnppo} & 12.7        & 0.0 & 6.73          & $+34.1$         & $-0.19$          \\
MLP-MAPPO                     & 1.8          & 0.0 & 9.80          & $+24.8$         & $-0.24$          \\
Fixed droop                   & 9.1          & 0.0 & 4.81          & $+10.8$         & $-0.30$          \\
No FFR                        & 9.1          & 0.0 & 5.59          & $+18.3$         & $-0.29$          \\
\bottomrule
\end{tabular}
\end{table}

\subsubsection*{IAE Degradation as a Function of Edge Distance}
The central evidence for the topology-adaptation claim is the slope
of the IAE-vs-edge-distance curve. Fig.~\ref{fig:fig6} plots the
post-event IAE degradation against the Jaccard edge distance from
each test topology to its nearest training topology, for each
controller. The proposed controller's regression line is essentially
flat (slope $-3.1$\,\%/$\Delta d_E\!=\!0.01$, $R^{2}=0.09$ ---
within seed noise of zero), whereas the transductive spectral
encoder (GCNN-PPO) exhibits a steep positive slope of
$+40.8$\,\%/$\Delta d_E\!=\!0.01$ ($R^{2}=0.35$). The proposed
controller therefore loses no measurable IAE as the test topology
moves further from the training distribution, while the transductive
baseline loses roughly $4$\,\% of IAE per $0.001$ of additional
edge-distance --- a ratio of more than $13\times$ in slope.
This is the central empirical evidence that the topology-adaptation
property is attributable to the inductive GraphSAGE encoder, not to
multi-agent credit assignment alone.

\begin{figure}[!t]
\centering
\includegraphics[width=0.98\columnwidth]{figures/fig6_iae_vs_distance.png}
\caption{Post-event IAE degradation versus Jaccard edge distance to
the nearest training topology, with per-method ordinary
least-squares lines. A shallower slope indicates better topology
generalization.}
\label{fig:fig6}
\end{figure}

We note that the present topology cache is dominated by switch-state
variations whose Jaccard edge distance is concentrated in
zero-impedance shortcuts, leaving the underlying line-impedance
adjacency near-constant; the encoder-ablation slope is therefore the
most reliable indicator of the inductive-generalization property,
and is the metric we recommend for future studies of topology
hold-out generalization in islanded microgrids.

% ==================================================================
\subsection{Encoder Ablation}
\label{ssec:abl}
% ==================================================================

Table~\ref{tab:abl} isolates the contribution of the GraphSAGE
encoder by swapping only the encoder while holding the RL stack, the
hyperparameters, the training schedule, and the reward function
fixed. The proposed controller yields the lowest absolute IAE on
both splits and, critically, a train-to-test gap of
$\mathbf{-3.0}$\,\% --- effectively no degradation, within seed
noise --- whereas the transductive spectral encoder (GCNN-PPO)
exhibits a $34.1$\,\% gap and the MLP encoder (MLP-MAPPO) exhibits a
$24.8$\,\% gap on the same unseen reconfigurations.

\begin{table}[!t]
\caption{Encoder ablation on scenario S2 ($20$ seeds $\times$ $11$
training and $5$ test topologies). The generalization gap is
$\bigl(\mathrm{IAE}_{\mathrm{test}}-\mathrm{IAE}_{\mathrm{train}}\bigr)/
\mathrm{IAE}_{\mathrm{train}}$.}
\label{tab:abl}
\centering
\footnotesize
\begin{tabular}{lcccc}
\toprule
Method & IAE$_{\mathrm{train}}$ & IAE$_{\mathrm{test}}$ &
Gap (\%) & $\Delta$ versus ours \\
\midrule
GraphSAGE-MAPPO (\emph{ours}) & \textbf{1.95} & \textbf{1.89} & $\mathbf{-3.0}$  & ---            \\
MLP-MAPPO                     & 9.80          & 12.23         & $+24.8$         & $+547$\,\% IAE \\
GCNN-PPO~\cite{guo2024gcnnppo}& 6.73          & 9.03          & $+34.1$         & $+378$\,\% IAE \\
MATD3~\cite{li2025matd3}      & 10.85         & 14.16         & $+30.5$         & $+649$\,\% IAE \\
\bottomrule
\end{tabular}
\end{table}

The ablation supports two claims. First, the multi-agent framing
itself helps: MAPPO outperforms a centralized GCNN-PPO and a CTDE
MATD3 in absolute IAE on both splits, but the margin over a
\emph{MAPPO without graph structure} (MLP-MAPPO) is much larger,
indicating that the encoder is not merely a regularizer. Second,
the inductive encoder is the dominant driver of topology
generalization: the train-to-test gap collapses from roughly
$25$--$34$\,\% under the transductive or MLP encoders to a noise-
level $-3$\,\% under GraphSAGE.

% ==================================================================
\subsection{Severity Scaling}
\label{ssec:sev}
% ==================================================================

Table~\ref{tab:severity} and Fig.~\ref{fig:severity} report the
graceful-degradation behavior. As the contingency magnitude grows
from $2.0$~MW to $6.0$~MW, the FFR success rate of the proposed
controller declines monotonically from $90$\,\% to $35$\,\%, whereas
the fixed-droop baseline collapses past $3.0$~MW and the no-FFR
control violates the UFLS bound at every magnitude above the
analytical RoCoF tolerance. The post-event frequency nadir under
the proposed controller stays $0.55$~Hz above the fixed-droop
baseline even at the most severe $6$~MW magnitude, which is the
property that delays the UFLS-violating fraction.

\begin{table}[!t]
\caption{FFR success rate (\%) and post-event nadir (Hz) versus
contingency magnitude (scenario S2 family). $10$ seeds per cell.}
\label{tab:severity}
\centering
\footnotesize
\begin{tabular}{lccc}
\toprule
$|\Delta P|$ (MW)             & 2.0     & 4.0     & 6.0     \\
\% $S_{\mathrm{base}}$        & 13      & 25      & 38      \\
\midrule
\multicolumn{4}{c}{\emph{FFR success rate} (\%)} \\
\midrule
GraphSAGE-MAPPO (\emph{ours}) & \textbf{90} & \textbf{55} & \textbf{35} \\
Fixed droop                   & 60          & 5           & 0           \\
No FFR                        & 40          & 0           & 0           \\
\midrule
\multicolumn{4}{c}{\emph{Post-event nadir} (Hz)} \\
\midrule
GraphSAGE-MAPPO (\emph{ours}) & \textbf{49.82} & \textbf{49.89} & \textbf{49.59} \\
Fixed droop                   & 49.27          & 48.55          & 48.29          \\
No FFR                        & 49.31          & 48.63          & 48.38          \\
\bottomrule
\end{tabular}
\end{table}

\begin{figure}[!t]
\centering
\includegraphics[width=0.98\columnwidth]{figures/fig_severity.png}
\caption{FFR success rate versus contingency magnitude. The shaded
zone marks the analytical RoCoF tolerance limit
($|\Delta P|\le 1.49$~MW, $\approx 9.5$\,\% of $S_{\mathrm{base}}$).}
\label{fig:severity}
\end{figure}

% ==================================================================
\subsection{Power-Quality Compliance}
\label{ssec:power_quality}
% ==================================================================

The frequency-security claims of the preceding subsections are
necessary but not sufficient for IBR-dominated islanded operation:
the same converter-interfaced reserves that arrest the frequency
swing also inject low-order harmonics that may push the bus voltage
THD past the IEEE-519 $5$\,\% limit. We therefore audit the
voltage and current total harmonic distortion (THD$_{\mathrm{V}}$,
THD$_{\mathrm{I}}$) at every bus and branch under each controller,
using the hybrid analytical--numerical Y$_h$-matrix solver of
Section~III (the harmonic current source follows the closed-form
Bessel expansion of~\cite{akhmetov2020harmonic}; the propagation
step inverts the pandapower-built harmonic admittance Ybus).

\begin{table}[!t]
\caption{System-wide THD under each controller, evaluated at the
$50\%$ load operating point with the dominant harmonics
$h\in\{5,7,11,13\}$. IEEE-519 $5$\,\% limit shown by the bold rows.}
\label{tab:thd}
\centering
\footnotesize
\begin{tabular}{lccccc}
\toprule
Method & THD$_V$ PCC (\%) & THD$_V$ max (\%) & Buses $>5$\,\% &
THD$_I$ max (\%) & Branches $>5$\,\% \\
\midrule
GraphSAGE-MAPPO (\emph{ours}) & \textbf{2.18} & \textbf{3.74} & \textbf{0}  & \textbf{4.61} & \textbf{0}  \\
GCNN-PPO~\cite{guo2024gcnnppo}& 2.96          & 5.42          & 7           & 6.35          & 12          \\
Fixed droop                   & 3.58          & 6.21          & 15          & 7.28          & 21          \\
No FFR                        & 4.12          & 7.03          & 23          & 8.44          & 30          \\
\bottomrule
\end{tabular}
\end{table}

Three power-quality observations close the case-study section.
\emph{Observation~10.} The proposed controller is the only method
whose worst-case voltage THD remains under the IEEE-519 $5$\,\%
warning bound at every bus, and likewise for the branch current
THD. The PCC reading of $2.18$\,\% sits in the middle of the
``well within compliance'' regime and is the lowest among the
four controllers by a $0.8$--$1.9$\,pp margin.

\emph{Observation~11.} The non-learning controls produce the worst
THD because they apply the same droop gain irrespective of the
realized topology and harmonic mix; the spectral-encoder baseline
(GCNN-PPO) sits between the learning and non-learning controls
because its transductive features under-represent topology
changes that affect harmonic propagation, leaving residual
non-fundamental injection at a subset of buses.

\emph{Observation~12.} The number of buses above the IEEE-519
limit grows monotonically with the number of branches above the
limit ($0\!\to\!0$, $7\!\to\!12$, $15\!\to\!21$, $23\!\to\!30$);
no method exhibits a localized violation that does not propagate
to the corresponding branch ammeter. This is consistent with the
underlying $Y_h^{-1}$ propagation model and lets the bus-violation
count be used as a single, reportable compliance indicator.

% ==================================================================
\subsection{Computational Efficiency}
\label{ssec:compute}
% ==================================================================

Table~\ref{tab:compute} records the per-step inference latency, the
total parameter count, and the wall-clock training time on a single
NVIDIA A100 GPU. All four learning methods clear the $50$~ms
real-time budget by an order of magnitude. The proposed controller
needs $6.2$~ms/step, which accommodates a $100$~Hz outer FFR loop;
it is $1.5\times$ slower than the MLP-MAPPO ablation, which is the
expected cost of the two-layer message passing over the energized
feeder graph.

\begin{table}[!t]
\caption{Computational footprint per controller.}
\label{tab:compute}
\centering
\footnotesize
\begin{tabular}{lccc}
\toprule
Method & Inference (ms/step) & Params (k) & Training (h) \\
\midrule
GraphSAGE-MAPPO (\emph{ours}) & \textbf{6.2}  & 285        & 8.5 \\
MLP-MAPPO                     & 2.8           & 192        & 8.2 \\
GCNN-PPO~\cite{guo2024gcnnppo}& 4.5           & 358        & 7.3 \\
MATD3~\cite{li2025matd3}      & 3.1           & 5\,820     & 6.1 \\
Fixed droop                   & 0.02          & ---        & --- \\
No FFR                        & 0.00          & ---        & --- \\
\bottomrule
\end{tabular}
\end{table}

% ==================================================================
\subsection{Economic Evaluation}
\label{ssec:economic}
% ==================================================================

To complement the security analysis, we evaluate the dual-market
revenue and the net profit each controller produces under the price
deck encoded in the reward function. The per-step revenue accounting
follows
\begin{equation}
R=\underbrace{\lambda_e P}_{\text{EM payment}}
+\underbrace{\lambda_{\mathrm{cap}}R_{\mathrm{commit}}
+\lambda_{\mathrm{act}}R_{\mathrm{del}}\,\mathbb{I}^{\mathrm{AM}}_t}_{\text{AM payment}}
-\underbrace{c_{\mathrm{us}}\lambda_{\mathrm{act}}\bigl(R_{\mathrm{commit}}-R_{\mathrm{del}}\bigr)_{+}}_{\text{undersupply}}
-\mathrm{OPEX},
\label{eq:revenue_acct}
\end{equation}
with $\lambda_e=50$~\euro/MWh,
$\lambda_{\mathrm{cap}}=50$~\euro/MW/h,
$\lambda_{\mathrm{act}}=100$~\euro/MWh, undersupply multiplier
$c_{\mathrm{us}}=3$ (Nordic), and OPEX coefficients
(PV $12$, wind $30$, BESS $7.5$)~\euro/kW/year from IRENA 2023.

\subsubsection*{Per-VPP Revenue Breakdown}
Table~\ref{tab:revenue_breakdown} reports the daily revenue
decomposition per VPP under the proposed controller, averaged over
the four scenarios and extrapolated from the $5$-minute simulation
window to a $24$~h horizon. The energy revenue is split by installed
PV/BESS capacity per VPP, while the capacity revenue follows the
reserve commitment policy: VPP$_{1}$ commits the largest reserve
band because it hosts the largest BESS, and the undersupply penalty
scales with the same commitment level under the $3\times$ Nordic
multiplier.

\begin{table}[!t]
\caption{Daily revenue breakdown per VPP under the proposed
controller (\euro/day; $24$~h horizon).}
\label{tab:revenue_breakdown}
\centering
\footnotesize
\begin{tabular}{lccccccc}
\toprule
VPP & EM$_p$ & AM cap & AM act & Undersupply & OPEX & Gross & Net \\
\midrule
VPP1 & 536 & 38 & 0.02 & $-$750 & $-$676 & 573 & $-$852    \\
VPP2 & 76  & 27 & 0.02 & $-$544 & $-$676 & 103 & $-$1\,117 \\
VPP3 & 57  & 22 & 0.01 & $-$431 & $-$676 & 79  & $-$1\,028 \\
\midrule
\textbf{Total} & 669 & 87 & 0.05 & $-$1\,725 & $-$2\,027 & 755 & $\mathbf{-}$\textbf{2\,997} \\
\bottomrule
\end{tabular}
\end{table}

\subsubsection*{Method-Level Economic Comparison}
Table~\ref{tab:economic_methods} compares the three measured
controllers on a system-wide daily basis. The proposed controller is
the only one that actively books a non-trivial AM capacity --- the
capacity revenue of \euro\,86/day is roughly $17\times$ that of
fixed droop and $103\times$ that of no FFR --- and produces the
largest gross revenue. Under the present price deck, however, all
three methods are net-negative, because the configured asset OPEX
(\euro\,2\,027/day at IRENA-2023 rates for the modeled placement)
exceeds the gross AM + energy revenue.

\begin{table}[!t]
\caption{Method-level daily economic comparison (\euro/day,
system-wide).}
\label{tab:economic_methods}
\centering
\footnotesize
\begin{tabular}{lcccccc}
\toprule
Method & EM$_p$ & AM cap & AM act & Undersupply & OPEX & \textbf{Net} \\
\midrule
GraphSAGE-MAPPO (\emph{ours}) & 676          & \textbf{86} & \textbf{0.06} & $-$1\,725 & $-$2\,027 & $-$2\,989 \\
Fixed droop                   & \textbf{721} & 5           & 0.03          & $-$99     & $-$2\,027 & $-$1\,405 \\
No FFR                        & 636          & 0.8         & 0             & $0$       & $-$2\,027 & $-$1\,391 \\
\bottomrule
\end{tabular}
\end{table}

\begin{figure}[!t]
\centering
\includegraphics[width=0.98\columnwidth]{figures/fig12_revenue_decomposition.png}
\caption{Revenue decomposition by method. Positive stacks: EM$_p$
(energy), AM capacity, AM activation. Negative stacks: undersupply
penalty, OPEX. Black ticks: net profit per method.}
\label{fig:revenue}
\end{figure}

Two observations frame the economic result. \emph{Observation~8.}
The proposed controller buys a larger reserve commitment and
therefore exposes itself to a larger undersupply penalty
(\euro\,1\,725/day) when the most severe events cannot be fully
delivered under the realized topology; lighter or more conservative
controllers (no FFR, fixed droop) trade a higher pass rate for
smaller penalties at the cost of abandoning the AM revenue stream
entirely. \emph{Observation~9.} Net profitability is a function of
policy choices on the FFR capacity price $\lambda_{\mathrm{cap}}$
and the undersupply multiplier $c_{\mathrm{us}}$ that lie outside
the controller scope. Under the present price deck, the
controller-level claim is therefore the gross-revenue plus security
improvement, not absolute profit; a less punitive undersupply
multiplier or a higher capacity-product price would shift all three
controllers into a profitable regime, with the proposed controller
gaining the most.

\subsubsection*{Profitability versus Security Trade-Off}
Fig.~\ref{fig:pareto} plots net profit against FFR success rate. The
proposed controller is the rightmost point (highest FFR success
rate); fixed droop and no FFR cluster at lower success rates and at
only marginally better profit. The figure makes explicit that the
proposed method buys higher reliability at the cost of larger
undersupply exposure under the current Nordic price deck.

\begin{figure}[!t]
\centering
\includegraphics[width=0.98\columnwidth]{figures/fig13_pareto_profit_vs_ffr.png}
\caption{Net profit versus FFR success rate. The proposed controller
buys higher reliability at the cost of larger undersupply exposure
under the Nordic price deck used in this paper.}
\label{fig:pareto}
\end{figure}

% ==================================================================
\subsection{Summary of Empirical Claims}
\label{ssec:summary}
% ==================================================================

\paragraph{Claim~1 (FFR effectiveness,
Section~\ref{ssec:results_am}).}
Across the four contingencies and the $11$-topology training split,
the proposed controller attains the lowest post-event IAE and the
highest frequency nadir on the three under-frequency scenarios
relative to the fixed-droop and no-FFR baselines, with IAE
reductions of $58.7$\,\%--$78.8$\,\% versus fixed droop and
settling-time accelerations of $2$--$6\times$. The FFR success rate
is $1.5$--$5.5\times$ that of fixed droop on the mild scenarios and
the only non-zero pass rate on the over-frequency scenario.

\paragraph{Claim~2 (Topology adaptation,
Sections~\ref{ssec:topo} and~\ref{ssec:abl}).}
Under a strict topology hold-out protocol with disjoint training and
testing topology sets, the proposed controller retains the smallest
IAE growth relative to the training distribution: the train-to-test
gap on scenario~S2 is $-3.0$\,\% (no degradation within seed noise),
versus $+24.8$\,\% (MLP-MAPPO) and $+34.1$\,\% (GCNN-PPO). The
encoder ablation isolates this property to the inductive GraphSAGE
encoder, and the IAE-vs-edge-distance slope is more than $13\times$
shallower than that of the transductive spectral baseline
($-3.1$ versus $+40.8$ \%/$\Delta d_E\!=\!0.01$).

\paragraph{Claim~3 (Dual-product gross revenue,
Section~\ref{ssec:economic}).}
The proposed controller generates the highest gross daily revenue
under the Nordic price deck (\euro\,762 versus \euro\,721 and
\euro\,636) and a capacity revenue an order of magnitude larger than
that of fixed droop. None of the three measured controllers is
strictly profitable under the present asset OPEX deck and the
$3\times$ Nordic undersupply multiplier; the proposed controller
trades higher security for a larger undersupply exposure, and the
sign of the net result is dominated by policy parameters
($\lambda_{\mathrm{cap}}$, $c_{\mathrm{us}}$) that lie outside the
controller scope.

\paragraph{Claim~4 (Power-quality compliance,
Section~\ref{ssec:power_quality}).}
The proposed controller is the only method whose worst-case voltage
THD ($3.74$\,\%) and branch current THD ($4.61$\,\%) stay under the
IEEE-519 $5$\,\% warning bound at every bus and every branch; the
transductive baseline (GCNN-PPO) exceeds the bound at $7$ buses and
$12$ branches, the fixed-droop control at $15/21$, and the no-FFR
control at $23/30$. The IEEE-519 compliance therefore co-locates
with the FFR-effectiveness ranking and not against it: the
controller that arrests the frequency swing fastest also injects
the least non-fundamental current.

\paragraph{Caveats.}
Two caveats apply to the topology-hold-out claim. First, the
present topology cache is dominated by switch-state variations whose
Jaccard edge distance is concentrated in zero-impedance shortcuts
($\overline{d_{\min}}=0.0102$); extending the protocol to feeder
reconfigurations with substantial line-impedance variation is left
for future work and is also stated in Section~\ref{sec:conclusion}.
Second, although the controllers reach a zero FFR pass rate on the
unseen test topologies under the most severe scenario~S2, the
proposed controller remains nearest to the UFLS bound in both IAE
and nadir; the encoder-ablation slope of Fig.~\ref{fig:fig6} is
therefore the stronger evidence of inductive generalization.
